% GENERATED FILE. Edit canonical lessons, then run npm run generate:books.
% canonical-source-hash: fnv1a64:185615d57429b471
% canonical-lessons: PA-C300-fasal, PA-C300-harhi, PA-C300-sauni, PA-C300-kapah, PA-C300-ganna

\chapter{Crops of the Year}
\label{ch:pa-a2-crops-of-the-year}
\hlchaptermodality{300}

\begin{chapteropening}
\textbf{By the end of this chapter:} \emph{I can name the crops of the farming year.}

Complete the last lesson of chapter 300: I can name the crops of the farming year.
\end{chapteropening}

\section[\texorpdfstring{\pa{ਫ਼ਸਲ} (fasal)}{ਫ਼ਸਲ (fasal)}]{\pa{ਫ਼ਸਲ} (fasal) — a crop; a harvest}

\label{lesson:PA-C300-fasal}



\begin{quote}
Before the new one: say the Punjabi for brackish, then the Punjabi for beautiful.
\end{quote}

\subsection*{You'll want to know: \pa{ਫ਼ਸਲ}}
\textbf{\pa{ਫ਼ਸਲ}} — \emph{fasal} — \textquotedblleft{}a crop; a harvest\textquotedblright{}.

\begin{quote}
\pa{ਅਗਲੇ} \pa{ਸਾਲ} \pa{ਫ਼ਸਲ} \pa{ਚੰਗੀ} \pa{ਹੋਵੇਗੀ।}
\end{quote}

\begin{quote}
\emph{agle sāl fasal chaṅgī hovegī.}
\end{quote}

\begin{quote}
Next year the crop will be good.
\end{quote}

A feminine noun.

\begin{cousinweb}
\textbf{\pa{ਫ਼ਸਲ}} (\emph{fasal}) is a loan, usually traced through Persian to Arabic \emph{fasl}, \textquotedblleft{}season\textquotedblright{}; its sense here, \textquotedblleft{}crop\textquotedblright{}, is a shift from that.
\end{cousinweb}

\begin{grammarlens}[title={What you've built}]
One more word for this chapter.
\end{grammarlens}

\subsection*{Your turn}
\begin{itemize}
\raggedright
  \item \emph{Say it:} \emph{fasal}
  \item \emph{Say it:} \emph{fasal}, once more
  \item \emph{Say it:} \emph{agle sāl fasal chaṅgī hovegī.} — Next year the crop will be good.
  \item \emph{Recall:} say the Punjabi for brackish, then the Punjabi for beautiful, then say \emph{fasal} again
\end{itemize}

\begin{tcolorbox}[breakable,colback=teal!4,colframe=teal!35!black,title={Before you move on}]
What does \pa{ਫ਼ਸਲ} mean? (\textquotedblleft{}A crop; a harvest\textquotedblright{}.) Say it once more.
\end{tcolorbox}
\section[\texorpdfstring{\pa{ਹਾੜ੍ਹੀ} (hāṛhī)}{ਹਾੜ੍ਹੀ (hāṛhī)}]{\pa{ਹਾੜ੍ਹੀ} (hāṛhī) — the spring harvest; the winter-sown crop}

\label{lesson:PA-C300-harhi}



\begin{quote}
Before the new one: say the Punjabi for beautiful, then the Punjabi for a crop.
\end{quote}

\subsection*{You'll want to know: \pa{ਹਾੜ੍ਹੀ}}
\textbf{\pa{ਹਾੜ੍ਹੀ}} — \emph{hāṛhī} — \textquotedblleft{}the spring harvest; the winter-sown crop\textquotedblright{}.

\begin{quote}
\pa{ਕਣਕ} \pa{ਹਾੜ੍ਹੀ} \pa{ਦੀ} \pa{ਫ਼ਸਲ} \pa{ਹੈ।}
\end{quote}

\begin{quote}
\emph{kaṇak hāṛhī dī fasal hai.}
\end{quote}

\begin{quote}
Wheat is a crop of the spring harvest.
\end{quote}

A feminine noun.

\begin{cousinweb}
\textbf{\pa{ਹਾੜ੍ਹੀ}} (\emph{hāṛhī}) is inherited from Sanskrit \emph{āṣāḍhīya}, \textquotedblleft{}pertaining to the month Āṣāḍha\textquotedblright{}, the month you know as \pa{ਹਾੜ}; CDIAL lists Punjabi \emph{hāṛhī} and \emph{hāṛī}, \textquotedblleft{}crops sown in winter\textquotedblright{}, under it (CDIAL 1474).
\end{cousinweb}

\begin{grammarlens}[title={What you've built}]
One more word for this chapter.
\end{grammarlens}

\subsection*{Your turn}
\begin{itemize}
\raggedright
  \item \emph{Say it:} \emph{hāṛhī}
  \item \emph{Say it:} \emph{hāṛhī}, once more
  \item \emph{Say it:} \emph{kaṇak hāṛhī dī fasal hai.} — Wheat is a crop of the spring harvest.
  \item \emph{Recall:} say the Punjabi for beautiful, then the Punjabi for a crop, then say \emph{hāṛhī} again
  \item \emph{Write it:} copy \textbf{\pa{ਹਾੜ੍ਹੀ}} once, checking the whole word
\end{itemize}

\begin{tcolorbox}[breakable,colback=teal!4,colframe=teal!35!black,title={Before you move on}]
What does \pa{ਹਾੜ੍ਹੀ} mean? (\textquotedblleft{}The spring harvest; the winter-sown crop\textquotedblright{}.) Say it once more.
\end{tcolorbox}
\section[\texorpdfstring{\pa{ਸਾਉਣੀ} (sāuṇī)}{ਸਾਉਣੀ (sāuṇī)}]{\pa{ਸਾਉਣੀ} (sāuṇī) — the autumn harvest; the monsoon-sown crop}

\label{lesson:PA-C300-sauni}



\begin{quote}
Before the new one: say the Punjabi for a crop, then the Punjabi for the spring harvest.
\end{quote}

\subsection*{You'll want to know: \pa{ਸਾਉਣੀ}}
\textbf{\pa{ਸਾਉਣੀ}} — \emph{sāuṇī} — \textquotedblleft{}the autumn harvest; the monsoon-sown crop\textquotedblright{}.

\begin{quote}
\pa{ਮੱਕੀ} \pa{ਸਾਉਣੀ} \pa{ਦੀ} \pa{ਫ਼ਸਲ} \pa{ਹੈ।}
\end{quote}

\begin{quote}
\emph{makkī sāuṇī dī fasal hai.}
\end{quote}

\begin{quote}
Maize is a crop of the autumn harvest.
\end{quote}

A feminine noun.

\begin{cousinweb}
\textbf{\pa{ਸਾਉਣੀ}} (\emph{sāuṇī}) is inherited from Sanskrit \emph{śrāvaṇika}, \textquotedblleft{}relating to the month Śrāvaṇa\textquotedblright{}; CDIAL lists Punjabi \emph{sāuṇī}, \textquotedblleft{}the autumn crop\textquotedblright{}, under it (CDIAL 12700).
\end{cousinweb}

\begin{grammarlens}[title={What you've built}]
One more word for this chapter.
\end{grammarlens}

\subsection*{Your turn}
\begin{itemize}
\raggedright
  \item \emph{Say it:} \emph{sāuṇī}
  \item \emph{Say it:} \emph{sāuṇī}, once more
  \item \emph{Say it:} \emph{makkī sāuṇī dī fasal hai.} — Maize is a crop of the autumn harvest.
  \item \emph{Recall:} say the Punjabi for a crop, then the Punjabi for the spring harvest, then say \emph{sāuṇī} again
  \item \emph{Write it:} copy \textbf{\pa{ਸਾਉਣੀ}} once, checking the whole word
\end{itemize}

\begin{tcolorbox}[breakable,colback=teal!4,colframe=teal!35!black,title={Before you move on}]
What does \pa{ਸਾਉਣੀ} mean? (\textquotedblleft{}The autumn harvest; the monsoon-sown crop\textquotedblright{}.) Say it once more.
\end{tcolorbox}
\section[\texorpdfstring{\pa{ਕਪਾਹ} (kapāh)}{ਕਪਾਹ (kapāh)}]{\pa{ਕਪਾਹ} (kapāh) — raw cotton; the cotton plant}

\label{lesson:PA-C300-kapah}



\begin{quote}
Before the new one: say the Punjabi for the spring harvest, then the Punjabi for the autumn harvest.
\end{quote}

\subsection*{You'll want to know: \pa{ਕਪਾਹ}}
\textbf{\pa{ਕਪਾਹ}} — \emph{kapāh} — \textquotedblleft{}raw cotton; the cotton plant\textquotedblright{}.

\begin{quote}
\pa{ਸਾਡੇ} \pa{ਖੇਤ} \pa{ਵਿੱਚ} \pa{ਕਪਾਹ} \pa{ਹੈ।}
\end{quote}

\begin{quote}
\emph{sāḍe khet vich kapāh hai.}
\end{quote}

\begin{quote}
There is cotton in our field.
\end{quote}

A feminine noun.

\begin{cousinweb}
\textbf{\pa{ਕਪਾਹ}} (\emph{kapāh}) is inherited from Sanskrit \emph{karpāsa}, \textquotedblleft{}the cotton plant\textquotedblright{}; CDIAL lists Punjabi \emph{kapāh}, \textquotedblleft{}cotton plant, raw cotton\textquotedblright{}, under it (CDIAL 2877).
\end{cousinweb}

\begin{grammarlens}[title={What you've built}]
One more word for this chapter.
\end{grammarlens}

\subsection*{Your turn}
\begin{itemize}
\raggedright
  \item \emph{Say it:} \emph{kapāh}
  \item \emph{Say it:} \emph{kapāh}, once more
  \item \emph{Say it:} \emph{sāḍe khet vich kapāh hai.} — There is cotton in our field.
  \item \emph{Recall:} say the Punjabi for the spring harvest, then the Punjabi for the autumn harvest, then say \emph{kapāh} again
  \item \emph{Write it:} copy \textbf{\pa{ਕਪਾਹ}} once, checking the whole word
\end{itemize}

\begin{tcolorbox}[breakable,colback=teal!4,colframe=teal!35!black,title={Before you move on}]
What does \pa{ਕਪਾਹ} mean? (\textquotedblleft{}Raw cotton; the cotton plant\textquotedblright{}.) Say it once more.
\end{tcolorbox}
\section[\texorpdfstring{\pa{ਗੰਨਾ} (gannā)}{ਗੰਨਾ (gannā)}]{\pa{ਗੰਨਾ} (gannā) — sugar cane}

\label{lesson:PA-C300-ganna}



\begin{quote}
Before the new one: say the Punjabi for the autumn harvest, then the Punjabi for raw cotton.
\end{quote}

\subsection*{You'll want to know: \pa{ਗੰਨਾ}}
\textbf{\pa{ਗੰਨਾ}} — \emph{gannā} — \textquotedblleft{}sugar cane\textquotedblright{}.

\begin{quote}
\pa{ਅਸੀਂ} \pa{ਗੰਨੇ} \pa{ਤੋਂ} \pa{ਗੁੜ} \pa{ਬਣਾਉਂਦੇ} \pa{ਹਾਂ।}
\end{quote}

\begin{quote}
\emph{asīṁ ganne ton guṛ baṇāunde hāṁ.}
\end{quote}

\begin{quote}
We make jaggery from sugar cane.
\end{quote}

A masculine noun.

\begin{cousinweb}
\textbf{\pa{ਗੰਨਾ}} (\emph{gannā}) is inherited from Sanskrit \emph{gaṇḍa}, \textquotedblleft{}joint of plant\textquotedblright{}; CDIAL lists Punjabi \emph{gannā}, \textquotedblleft{}sugarcane\textquotedblright{}, under it (CDIAL 3998).
\end{cousinweb}

\begin{grammarlens}[title={What you've built}]
One more word for this chapter.
\end{grammarlens}

\subsection*{Your turn}
\begin{itemize}
\raggedright
  \item \emph{Say it:} \emph{gannā}
  \item \emph{Say it:} \emph{gannā}, once more
  \item \emph{Say it:} \emph{asīṁ ganne ton guṛ baṇāunde hāṁ.} — We make jaggery from sugar cane.
  \item \emph{Recall:} say the Punjabi for the autumn harvest, then the Punjabi for raw cotton, then say \emph{gannā} again
  \item \emph{Write it:} copy \textbf{\pa{ਗੰਨਾ}} once, checking the whole word
\end{itemize}

\begin{tcolorbox}[breakable,colback=teal!4,colframe=teal!35!black,title={Before you move on}]
What does \pa{ਗੰਨਾ} mean? (\textquotedblleft{}Sugar cane\textquotedblright{}.) Say it once more.
\end{tcolorbox}
